% Generated draft: requires mathematical and source review.
\documentclass[11pt]{article}
\usepackage[margin=25mm]{geometry}
\usepackage{fontspec}
\setmainfont{Latin Modern Roman}
\usepackage{amsmath,amssymb,mathtools,unicode-math}
\setmathfont{Latin Modern Math}
\usepackage{enumitem,needspace,xurl,hyperref}
\setlength{\parindent}{0pt}
\setlength{\parskip}{5pt}
\setlength{\emergencystretch}{3em}
\newcommand{\N}{\mathbb{N}}
\newcommand{\Z}{\mathbb{Z}}
\newcommand{\Q}{\mathbb{Q}}
\newcommand{\R}{\mathbb{R}}
\newcommand{\C}{\mathbb{C}}
\newcommand{\F}{\mathbb{F}}
\newcommand{\A}{\mathbb{A}}
\newcommand{\PP}{\mathbb{P}}
\newcommand{\E}{\mathbb{E}}
\newcommand{\eps}{\varepsilon}
\newcommand{\dd}{\,\mathrm{d}}
\begin{document}
\section*{UnsolvedMath ID: 30002656}
\textbf{Source problem code:} catalogue entry
\subsection{Definitions and mathematical statement}
Let \((X,\Delta)\) denote a log canonical pair of dimension \(n\). In the usual discrepancy notation, if \(f:W\to X\) is a birational model on which a prime divisor \(F\) appears, its log discrepancy is
\[
 a(F;X,\Delta)=1+\operatorname{coeff}_{F}\bigl(K_W-f^*(K_X+\Delta)\bigr).
\]
This notation requires the usual hypotheses permitting the pullback of \(K_X+\Delta\). Log canonicity means that all such log discrepancies are nonnegative. A non-klt centre of a log canonical pair is the centre of a divisor with log discrepancy zero.

For a projective birational morphism \(\pi:Y\to X\), write \(\widetilde\Delta\) for the strict transform of \(\Delta\), and write
\[
 \Gamma=\widetilde\Delta+E.
\]
The conventional interpretation of \(E\) is the reduced sum of the \(\pi\)-exceptional prime divisors on \(Y\); the supplied record does not explicitly define \(E\), so this interpretation requires confirmation.

\textbf{Numerical existence question, with an unresolved local condition.}
For every positive integer \(n\), does there exist a nonnegative integer \(s=s(n)\) such that every \(n\)-dimensional log canonical pair \((X,\Delta)\) admits a projective birational morphism \(\pi:Y\to X\) satisfying
\[
 \#\{F\subset Y:F\text{ is a }\pi\text{-exceptional prime divisor}\}\leq s(n),
 \qquad
 a(F;X,\Delta)=0
 \quad\text{for every such }F,
\]
and the prescribed local condition on \((Y,\Gamma)\) near the generic point of every non-klt centre?

The supplied formulations prescribe different local conditions:
\begin{itemize}
\item The original statement requires \((Y,\Gamma)\) to be \emph{toroidal} near these generic points.
\item The clean statement and current statement require \((Y,\Gamma)\) to be \emph{quotient-dlt} near these generic points.
\end{itemize}
These are retained as distinct candidate formulations; their equivalence is not asserted. The record does not explicitly identify whether the non-klt centres belong to \((X,\Delta)\) or \((Y,\Gamma)\), nor specify how a neighbourhood on \(Y\) is to be understood if the centres are on \(X\). It also leaves the ground field, characteristic, coefficient domain of \(\Delta\), and precise local meanings of toroidal and quotient-dlt unspecified. Thus the supplied evidence does not determine a single fully specified conjecture.

The supplied literature assessment further describes the extracted assertion as only a numerical interpretation of an economical-resolution conjecture whose full statement is omitted. The question above records the supplied numerical assertion and does not reconstruct that missing conjecture.
\subsection{Short English statement}
The proposed bound depends only on dimension: every log canonical pair should admit a projective birational modification extracting at most that many divisors, all with log discrepancy zero, and satisfying a local condition near non-klt centres. The original record asks for toroidal structure, whereas the cleaned version asks for quotient-dlt structure. This difference, the ambient pair of the centres, and missing foundational hypotheses require review. The supplied open label is dated curation, not an independently verified status.
\subsection{Sources}
\begin{itemize}
\item Local UnsolvedMath registry, \texttt{data/unsolvedmath/problems.json}, permanent ID \(30002656\), source problem code \texttt{catalogue entry}; supplied title: \emph{Economical Toroidal Resolutions of Log Canonical Pairs}. The record belongs to the Oberwolfach Reports Open Problems collection. No individual proposer is supplied; the proposed year is recorded as 2014.
\item Supplied report citation: \emph{Komplexe Analysis (2015)}, Oberwolfach Reports, pp.\ 2177--2234, DOI \texttt{10.4171/owr/2014/39}. Source URL and supplied literature-source link labelled ``original Oberwolfach report'': \url{https://doi.org/10.4171/owr/2014/39}. The background describes these reports as records of mathematical workshops and their open-problem discussions. The bibliographic details and report contents have not been independently checked here.
\item Supplied literature assessment, checked at the exact recorded date \texttt{2026-08-21}: status labelled ``open''; the assessment says the extracted sentence gives only a numerical interpretation of an economical-resolution conjecture and omits the conjecture itself, making the claim unclear and leading to its normalization to open. This is supplied dated curation derived from the OWR record, not an independently verified literature assessment; time-sensitive status and cited claims require independent verification.
\item Supplied statement-verification metadata labels the statement \texttt{corrected\_verified} and reports comparison with the official Oberwolfach report, claiming that the clean statement restores source hypotheses and notation. This is a supplied verification claim, not verification performed for this draft; the toroidal versus quotient-dlt discrepancy remains unresolved.
\end{itemize}
\end{document}
